\documentclass[11pt]{article}

\usepackage[margin=1in]{geometry}
\usepackage{amsmath,amssymb,amsthm,mathtools}
\usepackage{enumitem}
\usepackage[hidelinks,hypertexnames=false]{hyperref}
\usepackage{xcolor}

\newtheorem{assumption}{Assumption}[section]
\newtheorem{theorem}[assumption]{Theorem}
\newtheorem{proposition}[assumption]{Proposition}
\newtheorem{lemma}[assumption]{Lemma}
\newtheorem{corollary}[assumption]{Corollary}
\newtheorem{definition}[assumption]{Definition}
\theoremstyle{remark}
\newtheorem{remark}[assumption]{Remark}

\newcommand{\R}{\mathbb R}
\newcommand{\E}{\mathbb E}
\newcommand{\ind}{\mathbf 1}
\newcommand{\dd}{\,d\nu}
\newcommand{\cA}{\mathcal A}
\newcommand{\cY}{\mathcal Y}
\newcommand{\cQ}{\mathcal Q}
\DeclareMathOperator{\Capacity}{Cap}
\DeclareMathOperator{\Curv}{Curv}
\DeclareMathOperator{\Width}{Width}
\newcommand{\ubar}[1]{\underline{#1}}
\newcommand{\proofsketch}{\paragraph{Proof sketch.}}
\newcommand{\todo}[1]{\par\smallskip\noindent\textcolor{red!70!black}{\textbf{Open point.} #1}\par\smallskip}

\title{Mathematical Plan for the Generalized Paper}
\author{Azevedo--Wolff (working document)}
\date{\today}

\begin{document}
\maketitle

\begin{abstract}
This is the master mathematical document for the generalized version of the paper.
It fixes notation, records the assumptions, and states the results in their intended
logical order.  A proof sketch follows every non-obvious statement.  The document
covers the relaxed solution, the finite-$\lambda$ concavity argument, high feasible
reservation utility, utility and distribution examples, and the appendix-only
asymptotic-affinity extension.  It is a planning document rather than manuscript
prose; theorem names and labels are intended to survive into the paper when possible.
\end{abstract}

\tableofcontents

\section{Primitives and contract representations}

\subsection{Actions, output, and scores}

\begin{definition}[Actions and output]
The agent chooses an action $a$ in a compact interval $\cA\subseteq\R_+$.
The intended action $a_0$ is in the interior of $\cA$.  Output has common
support $\cY\subseteq\R$ for every $a\in\cA$ and has density or probability mass
function $f(\cdot\mid a)>0$ on $\cY$ with respect to a common dominating measure
$\nu$.  Thus
\[
  \E_a[h(Y)]=\int_{\cY}h(y)f(y\mid a)\dd(y).
\]
Counting measure is used for discrete output.
\end{definition}

\begin{assumption}[Distributional regularity]
\label{ass:distribution-regularity}
For every $y\in\cY$, the map $a\mapsto f(y\mid a)$ is twice continuously
differentiable on an open set containing $\cA$.  For every $a^*\in\cA$, there are a relative neighborhood
$N_{a^*}$ of $a^*$ in $\cA$ and a nonnegative $\nu$-integrable function
$g_{a^*}$ such that
\begin{equation}
 \sup_{a\in N_{a^*}}
 \left\{
 |f(y\mid a)|+|f_a(y\mid a)|+|f_{aa}(y\mid a)|
 \right\}
 \le g_{a^*}(y)
 \quad\text{for $\nu$-almost every }y.
 \tag{DensDom}\label{eq:density-domination}
\end{equation}
\end{assumption}

Assumption~\ref{ass:distribution-regularity} and differentiation under the
integral sign imply, for every $a\in\cA$,
\[
 \int f_a(y\mid a)\dd(y)=0,
 \qquad
 \int f_{aa}(y\mid a)\dd(y)=0.
\]

\begin{definition}[Twice differentiable integrand class]
\label{def:H2}
Let $\mathcal H_f^2$ be the set of measurable functions $h:\cY\to\R$ such that,
for every $a^*\in\cA$, there are a relative neighborhood $N_{a^*}$ and a
nonnegative $\nu$-integrable function $g_{h,a^*}$ satisfying
\begin{equation}
 \sup_{a\in N_{a^*}}
 |h(y)|
 \left\{
 |f(y\mid a)|+|f_a(y\mid a)|+|f_{aa}(y\mid a)|
 \right\}
 \le g_{h,a^*}(y)
 \quad\text{for $\nu$-almost every }y.
 \tag{$H^2$}\label{eq:H2}
\end{equation}
For every $h\in\mathcal H_f^2$, the map
$a\mapsto\int h(y)f(y\mid a)\dd(y)$ is twice continuously differentiable and
its first two derivatives are obtained by differentiating $f$ under the integral.
\end{definition}

\begin{definition}[Score]
The score is
\[
 S(y\mid a):=\partial_a\log f(y\mid a)=\frac{f_a(y\mid a)}{f(y\mid a)}.
\]
After fixing $a_0$, write
\[
 s(y):=S(y\mid a_0),\qquad f_0(y):=f(y\mid a_0).
\]
Assumption~\ref{ass:distribution-regularity} gives
\[
 \E_{a_0}[s(Y)]=\int s(y)f_0(y)\dd(y)=0.
\]
\end{definition}

No monotone likelihood ratio property, interval support, or invertibility of the
score is maintained in the general results.

\subsection{Utility, wealth, payment, and cost}

\begin{definition}[Wealth and payment]
The agent has initial wealth $w_0$.  A payment contract is a measurable function
$w:\cY\to\R_+$.  Final wealth is
\[
 W(y):=w_0+w(y),
\]
so limited liability is equivalently $W(y)\ge w_0$.
\end{definition}

The admissible wealth domain is utility-specific.  Log and CRRA utility require
$w_0>0$; CARA utility permits $w_0\in\R$.

\begin{assumption}[Utility and effort cost]
\label{ass:utility-cost}
Utility $u$ is defined on an interval containing $[w_0,\infty)$ and is smooth,
strictly increasing, and strictly concave.  Moreover,
\[
 \lim_{W\to\infty}u'(W)=0,
 \qquad
 \bar u:=\lim_{W\to\infty}u(W)\in\R\cup\{\infty\}.
\]
Define
\[
 \underline u:=u(w_0),
 \qquad
 \Delta u:=\bar u-\underline u\in(0,\infty].
\]
The effort cost $c$ is smooth and strictly convex, and $c'(a_0)>0$.
\end{assumption}

Uniform concavity results will additionally impose
\begin{equation}
 \underline c'':=\inf_{a\in\cA}c''(a)>0.
 \tag{CostCurv}\label{eq:cost-curv}
\end{equation}
This is not needed to characterize the relaxed solution.

\begin{definition}[Three contract representations]
A payment contract $w$ has associated final-wealth and utility contracts
\[
 W(y)=w_0+w(y),
 \qquad
 v(y)=u(W(y)).
\]
Let $k:=u^{-1}$ on $[\underline u,\bar u)$.  Conversely,
\[
 W(y)=k(v(y)),
 \qquad
 w(y)=k(v(y))-w_0.
\]
The agent's expected utility and expected compensation cost are
\begin{align*}
 U(v,a)&:=\int v(y)f(y\mid a)\dd(y)-c(a),\\
 C(v,a)&:=\int [k(v(y))-w_0]f(y\mid a)\dd(y).
\end{align*}
Formally, the first argument of $U$ and $C$ is always a utility contract.  We do
not overload these functionals by writing $U(w,a)$ or $U(W,a)$; payment and
final-wealth contracts are first transformed into utility units.  We use $C$,
rather than $W(v,a)$, for expected compensation cost so that $W(y)$ remains
available for final wealth.
\end{definition}

\subsection{Link functions and canonical contracts}

\begin{definition}[Marginal-cost floor and links]
The marginal cost of utility at the limited-liability floor is
\[
 m_0:=k'(\underline u)=\frac{1}{u'(w_0)}.
\]
On the nonbinding region define the wealth and utility links
\[
 L_0(z):=(u')^{-1}(1/z),
 \qquad
 \ell_0(z):=u(L_0(z))=(k')^{-1}(z).
\]
The limited-liability links are
\[
 L(z):=L_0(z\vee m_0),
 \qquad
 \ell(z):=\ell_0(z\vee m_0)=u(L(z)).
\]
Thus $L(z)=w_0$ and $\ell(z)=\underline u$ for $z\le m_0$.
At $m_0$, derivatives always mean the relevant one-sided derivative.
\end{definition}

On the nonbinding region,
\[
 u'(L_0(z))=\frac1z,
 \qquad
 \ell_0'(z)=\frac{L_0'(z)}{z}.
\]

\begin{definition}[Canonical contracts]
For multipliers $\lambda,\mu\in\R$, define
\begin{align*}
 W_{\lambda,\mu}(y)&:=L(\lambda+\mu s(y)),\\
 w_{\lambda,\mu}(y)&:=L(\lambda+\mu s(y))-w_0,\\
 v_{\lambda,\mu}(y)&:=\ell(\lambda+\mu s(y)).
\end{align*}
The first is final wealth, the second is payment, and the third is utility.
Subscripts, rather than a vertical bar, distinguish contract parameters from the
conditioning notation in $f(y\mid a)$.
\end{definition}

\begin{assumption}[Regularity of feasible contracts]
\label{ass:contract-regularity}
The feasible utility-contract set $\mathcal V$ is convex.  Every
$v\in\mathcal V$ takes values in $[\underline u,\bar u)$ and belongs to
$\mathcal H_f^2$.  Every canonical contract $v_{\lambda,\mu}$ belongs to
$\mathcal V$ and has finite expected compensation cost at each $a\in\cA$.
Moreover, for every compact $K\subset\R^2$, there is a nonnegative
$\nu$-integrable function $G_K$ such that, for every $(\lambda,\mu)\in K$,
\begin{multline}
 \bigl|\ell(\lambda+\mu s(y))\bigr|
 \bigl(1+|s(y)|\bigr)f_0(y)
 \\
 {}+\bigl[L(\lambda+\mu s(y))-w_0\bigr]f_0(y)
 \le G_K(y)
 \quad\text{for $\nu$-almost every }y.
 \tag{CanDom}\label{eq:canonical-domination}
\end{multline}
\end{assumption}

Condition~\eqref{eq:H2} is the exact domination condition used to differentiate
agent utility with respect to action.  Condition~\eqref{eq:canonical-domination}
is the local uniform domination used for continuity in $(\lambda,\mu)$ of the
incentive functional, expected utility, and expected compensation cost.

\section{The relaxed problem}

The relaxed cost-minimization problem for $(a_0,\bar U)$ minimizes
$C(v,a_0)$ subject to
\begin{align}
 U(v,a_0)&\ge \bar U, \tag{IR}\label{eq:IR}\\
 U_a(v,a_0)&=0. \tag{LIC}\label{eq:LIC}
\end{align}
Under Assumption~\ref{ass:contract-regularity}, LIC is
\[
 \int v(y)s(y)f_0(y)\dd(y)=c'(a_0).
\]

\subsection{Local implementability and the LIC path}

\begin{assumption}[Local implementability]
\label{ass:implementability}
The intended action satisfies
\begin{equation}
 \Delta u\int_{\{s>0\}}s(y)f_0(y)\dd(y)>c'(a_0).
 \tag{Imp}\label{eq:implementability}
\end{equation}
When $\Delta u=\infty$, this means that the positive score has positive
first moment.
\end{assumption}

\begin{definition}[Incentive functional]
For fixed $\lambda$ and $\mu$, let
\[
 I(\lambda,\mu):=
 \int \ell(\lambda+\mu s(y))s(y)f_0(y)\dd(y).
\]
Thus a canonical contract satisfies LIC exactly when
$I(\lambda,\mu)=c'(a_0)$.
\end{definition}

\begin{lemma}[Unique LIC multiplier]
\label{lem:unique-mu}
Suppose Assumptions~\ref{ass:distribution-regularity}--\ref{ass:implementability}
hold.  For every
$\lambda\ge0$, there is a unique $\widetilde\mu(\lambda)>0$ such that
\[
 I(\lambda,\widetilde\mu(\lambda))=c'(a_0).
\]
The map $\lambda\mapsto\widetilde\mu(\lambda)$ is continuous.
\end{lemma}

\proofsketch
Because $\E_{a_0}[s]=0$, $I(\lambda,0)=0$.  Monotonicity of $\ell$ implies that
$I$ is weakly increasing in $\mu$; where differentiation is justified,
\[
 I_\mu(\lambda,\mu)=
 \int \ell'(\lambda+\mu s(y))s(y)^2f_0(y)\dd(y)\ge0.
\]
After subtracting the constant $\ell(\lambda)$ and using $\E_{a_0}[s]=0$,
the integrand is nonnegative and increases pointwise in $\mu$.  Hence monotone
convergence applies.  As $\mu\to\infty$,
\[
 \ell(\lambda+\mu s(y))\longrightarrow
 \begin{cases}
  \bar u,&s(y)>0,\\
  \ell(\lambda),&s(y)=0,\\
  \underline u,&s(y)<0.
 \end{cases}
\]
Consequently, with an infinite limit allowed,
\[
 I(\lambda,\mu)\longrightarrow
 \Delta u\int_{\{s>0\}}s(y)f_0(y)\dd(y)>c'(a_0).
\]
Existence follows by continuity.  At any LIC solution, positive generated
incentives imply that $\ell'$ is positive on a set with $s\ne0$, yielding strict
monotonicity at the solution and hence uniqueness.  Continuity follows from the
same strict crossing argument.

\begin{definition}[LIC path]
Define the utility contract
\[
 v_\lambda:=v_{\lambda,\widetilde\mu(\lambda)},
 \qquad
 v_\lambda(y)=\ell(\lambda+\widetilde\mu(\lambda)s(y)),
\]
with associated expected utility and cost
\begin{align*}
 \widetilde U(\lambda)&:=U(v_\lambda,a_0),\\
 \widetilde C(\lambda)&:=C(v_\lambda,a_0).
\end{align*}
\end{definition}

\subsection{Lagrangian and Pareto characterization}

\begin{lemma}[Pointwise Lagrangian minimization]
\label{lem:pointwise}
For multipliers $\lambda\ge0$ and $\mu\in\R$, consider
\[
 \mathcal L(v,\lambda,\mu)
 :=C(v,a_0)+\lambda[\bar U-U(v,a_0)]-\mu U_a(v,a_0).
\]
Among feasible utility contracts, its almost-everywhere unique minimizer is
$v_{\lambda,\mu}$.
\end{lemma}

\proofsketch
After dropping terms independent of $v$, the integrand is
\[
 k(v(y))-w_0-[\lambda+\mu s(y)]v(y).
\]
Strict convexity of $k$ gives a unique pointwise minimizer.  The interior first-order
condition is $k'(v)=\lambda+\mu s$; imposing $v\ge\underline u$ replaces the
right-hand side by its maximum with $m_0$, which gives $v=\ell(\lambda+\mu s)$.

\begin{lemma}[Relaxed Pareto path]
\label{lem:pareto-path}
For every $\lambda\ge0$, $v_\lambda$ is the almost-everywhere
unique minimizer of
\[
 C(v,a_0)-\lambda U(v,a_0)
\]
among contracts satisfying LIC.  The function $\widetilde U$ is continuous and
strictly increasing.
\end{lemma}

\proofsketch
The first claim follows from Lemma~\ref{lem:pointwise} with
$\mu=\widetilde\mu(\lambda)$.  For $\lambda_1<\lambda_2$, the two Pareto inequalities
imply
\[
 (\lambda_2-\lambda_1)
 [\widetilde U(\lambda_2)-\widetilde U(\lambda_1)]\ge0.
\]
If equality held, strict convexity of $k$ and uniqueness would force the two
contracts to agree almost everywhere, contradicting their distinct supporting
multipliers under score nondegeneracy.  Continuity follows from continuity of the
canonical path and dominated convergence.

\begin{definition}[Reservation-utility frontiers]
Define
\[
 U_L:=\widetilde U(0),
 \qquad
 U_R:=\lim_{\lambda\to\infty}\widetilde U(\lambda)
 \in\R\cup\{\infty\}.
\]
The limit exists because $\widetilde U$ is increasing.  The phrase
\emph{high feasible reservation utility} means
\[
 \bar U\uparrow U_R.
\]
If $\bar u=\infty$, then $U_R=\infty$.  If $\bar u<\infty$, then
$U_R\le \bar u-c(a_0)<\infty$.
\end{definition}

\begin{proposition}[Generalized relaxed solution]
\label{prop:relaxed-solution}
Suppose Assumptions~\ref{ass:distribution-regularity}--\ref{ass:implementability}
hold and the preceding nondegeneracy conditions hold.  Then:
\begin{enumerate}[label=(\roman*)]
 \item for every $\bar U<U_R$, the relaxed problem has an almost-everywhere
 unique solution;
 \item there are multipliers $\lambda^*(\bar U)\ge0$ and
 $\mu^*(\bar U)=\widetilde\mu(\lambda^*(\bar U))>0$ such that
 \[
  v^*_{\bar U}=v_{\lambda^*(\bar U),\mu^*(\bar U)},
  \qquad
  v^*_{\bar U}(y)
  =\ell\bigl(\lambda^*(\bar U)+\mu^*(\bar U)s(y)\bigr);
 \]
 \item if $\bar U\le U_L$, then $\lambda^*(\bar U)=0$ and both the relaxed
 contract and its expected cost are constant in $\bar U$;
 \item on $(U_L,U_R)$, $\lambda^*(\bar U)$ is strictly increasing and
 \[
  \lambda^*(\bar U)\longrightarrow\infty
  \quad\text{as}\quad \bar U\uparrow U_R;
 \]
 \item the relaxed value $C^*(\bar U)$ is strictly increasing and strictly convex
 on $(U_L,U_R)$.
\end{enumerate}
\end{proposition}

\proofsketch
For $\bar U>U_L$, strict continuity of $\widetilde U$ identifies a unique
$\lambda$ satisfying $\widetilde U(\lambda)=\bar U$.  The Pareto-path lemma and
the supporting-multiplier argument prove optimality and uniqueness.  Below
$U_L$, the IR constraint is slack and the solution is the $\lambda=0$
contract.  The inverse relationship between $\widetilde U$ and $\lambda^*$ gives
the frontier limit.  Strict convexity of the value follows from distinct supporting
slopes and strict convexity of compensation cost.

\todo{Before manuscript drafting, make the functional-analytic existence assumptions
fully explicit.  In particular, verify that bounded utility does not require a
separate closure argument at the unattained pointwise upper bound $\bar u$.}

\section{Score-tail geometry and the limited-liability kink}

\begin{assumption}[Nonbinding-link concavity]
\label{ass:link-concavity}
The unconstrained utility link $\ell_0$ is nondecreasing and concave on
$(m_0,\infty)$.  The limited-liability link $\ell(z)=\ell_0(z\vee m_0)$ need not
be globally concave because its slope jumps upward at $m_0$.
\end{assumption}

\begin{definition}[Limited-liability score cutoff]
For $\mu>0$, define
\[
 q^{\mathrm{LL}}(\lambda,\mu):=\frac{m_0-\lambda}{\mu}.
\]
Limited liability binds precisely on
\[
 \{y:s(y)\le q^{\mathrm{LL}}(\lambda,\mu)\}.
\]
Along the LIC path write
\[
 \widetilde q^{\mathrm{LL}}(\lambda)
 :=q^{\mathrm{LL}}(\lambda,\widetilde\mu(\lambda)).
\]
\end{definition}

\begin{definition}[Tail-safe cutoffs]
A number $q\le0$ is \emph{tail-safe} if
\[
 f_{aa}(y\mid a)\ge0
 \quad\text{for every }a\in\cA
 \text{ and every }y\text{ with }s(y)\le q.
 \tag{TS}\label{eq:tail-safe}
\]
Let
\[
 \cQ_{\mathrm{safe}}:=\{q\le0:q\text{ is tail-safe}\}.
\]
Define incentive capacity and paid-region positive curvature by
\begin{align*}
 \Capacity(q)&:=-\int_{\{s(y)\le q\}}s(y)f_0(y)\dd(y),\\
 \Curv(q)&:=\sup_{a\in\cA}
 \int_{\{s(y)>q\}}[s(y)f_{aa}(y\mid a)]_+\dd(y).
\end{align*}
The set $\cQ_{\mathrm{safe}}$ is a lower interval.  Define its right endpoint,
safe incentive capacity, and best paid-region curvature bound by
\begin{align*}
 \bar q_{\mathrm{safe}}&:=\sup\cQ_{\mathrm{safe}},\\
 \Capacity_{\mathrm{safe}}
 &:=\sup_{q<\bar q_{\mathrm{safe}}}\Capacity(q)
 =-\int_{\{s(y)<\bar q_{\mathrm{safe}}\}}s(y)f_0(y)\dd(y),\\
 \Curv_{\mathrm{safe}}
 &:=\inf_{q\in\cQ_{\mathrm{safe}}}\Curv(q).
\end{align*}
The supremum is interpreted as zero if there are no interior safe cutoffs, and
the infimum is interpreted as infinity if there are no safe cutoffs.  For the
logarithmic result, also define
\[
 \underline s:=\operatorname*{ess\,inf}_{Y\sim f_0}s(Y),
 \qquad
 \Width_{\mathrm{safe}}:=\bar q_{\mathrm{safe}}-\underline s.
\]
\end{definition}

Throughout, when $\Delta u=\infty$, a capacity condition of the form
$\Delta u\,K>c'(a_0)$ means $K>0$.  This convention applies both to
$K=\Capacity(q)$ and to $K=\Capacity_{\mathrm{safe}}$.

The weak inequality in $\{s(y)\le q\}$ assigns a score atom at $q$ to the safe,
no-payment side.  Its strict complement appears in $\Curv(q)$.  The function
$\Capacity(q)$ is increasing in $q$, whereas $\Curv(q)$ is decreasing.  Safe
incentive capacity excludes any capacity concentrated at a rightmost safe score
atom with no larger safe cutoff.

\section{Finite-\texorpdfstring{$\lambda$}{lambda} concavity lemmas}

\begin{definition}[Trial multiplier]
For $\lambda>m_0$ and $q<0$, let
\[
 \mu_q(\lambda):=\frac{m_0-\lambda}{q}>0.
\]
This multiplier places the kink at $q$:
$\lambda+\mu_q(\lambda)q=m_0$.
\end{definition}

\begin{lemma}[Kink placement]
\label{lem:kink-placement}
Fix $q<0$.  If
\[
 \Delta u\,\Capacity(q)>c'(a_0),
\]
then, for all sufficiently large $\lambda$,
\[
 \widetilde q^{\mathrm{LL}}(\lambda)\le q,
 \qquad
 \widetilde\mu(\lambda)\le\mu_q(\lambda)=O(\lambda).
\]
Tail safety of $q$ is not needed for this lemma, although later applications use
a tail-safe $q$.
\end{lemma}

\proofsketch
At $\mu_q(\lambda)$, scores below $q$ receive $\underline u$.  Subtracting the
constant $\ell(\lambda)$ and using $\E_{a_0}[s]=0$ gives
\[
 I(\lambda,\mu_q(\lambda))
 \ge [\ell(\lambda)-\underline u]\Capacity(q).
\]
All omitted terms are nonnegative: for negative scores above $q$, both the utility
difference and the score are nonpositive; for positive scores both are nonnegative.
The strict capacity condition makes the right side exceed $c'(a_0)$ eventually.
Monotonicity of $I$ then gives $\widetilde\mu\le\mu_q$.  Since $m_0-\lambda<0$,
dividing by the smaller positive multiplier moves the cutoff weakly left:
$\widetilde q^{\mathrm{LL}}(\lambda)\le q$.

\begin{definition}[Secant slope]
Along the LIC path, abbreviate $\widetilde\mu=\widetilde\mu(\lambda)$ and define
\[
 D_\lambda\ell(y):=
 \frac{\ell(\lambda+\widetilde\mu s(y))-\ell(\lambda)}{\widetilde\mu s(y)}
\]
when $s(y)\ne0$, with its continuous extension when $s(y)=0$.
\end{definition}

\begin{lemma}[Curvature identity]
\label{lem:curvature-identity}
For every $a\in\cA$,
\begin{equation}
 U_{aa}(v_\lambda,a)
 =\widetilde\mu
 \int D_\lambda\ell(y)s(y)f_{aa}(y\mid a)\dd(y)-c''(a).
 \tag{Curv}\label{eq:curvature-identity}
\end{equation}
\end{lemma}

\proofsketch
Add and subtract the constant $\ell(\lambda)$ from the utility contract and use
$\int f_{aa}(y\mid a)\dd(y)=0$.  Differentiation under the integral sign then
gives the result.

\begin{lemma}[Curvature split]
\label{lem:curvature-split}
Fix $q\in\cQ_{\mathrm{safe}}$ with $\Curv(q)<\infty$.  If the limited-liability cutoff
is weakly below $q$,
\[
 \widetilde q^{\mathrm{LL}}(\lambda)\le q,
\]
then, for every $a\in\cA$,
\begin{equation}
 U_{aa}(v_\lambda,a)
 \le
 \widetilde\mu\,\ell_0'(\lambda+\widetilde\mu q)\Curv(q)-c''(a).
 \tag{Split}\label{eq:curvature-split}
\end{equation}
If the argument equals $m_0$, use the right derivative.
\end{lemma}

\proofsketch
On $\{s(y)\le q\}$, tail safety gives $f_{aa}\ge0$, while $s\le0$ and monotonicity gives
$D_\lambda\ell\ge0$.  Hence this region contributes weakly negatively to
\eqref{eq:curvature-identity}.  On $\{s>q\}$ the contract is nonbinding.
Concavity of $\ell_0$ bounds every relevant secant slope by
$\ell_0'(\lambda+\widetilde\mu q)$.  Retaining only the positive part of
$s f_{aa}$ and taking the supremum over actions gives \eqref{eq:curvature-split}.

\begin{proposition}[Deterministic finite-$\lambda$ criterion]
\label{prop:deterministic-criterion}
Fix $q_0<q_1\le0$ with $q_0,q_1\in\cQ_{\mathrm{safe}}$.  Suppose $q_0$
has sufficient incentive capacity and paid-region curvature is finite at $q_1$:
\[
 \Delta u\,\Capacity(q_0)>c'(a_0),
 \qquad
 \Curv(q_1)<\infty.
\]
For all sufficiently large $\lambda$,
\begin{equation}
 \sup_{a\in\cA}U_{aa}(v_\lambda,a)
 \le
 \mu_{q_0}(\lambda)
 \ell_0'(\lambda+\mu_{q_0}(\lambda)q_1)\Curv(q_1)-\underline c''.
 \tag{DC-$\lambda$}\label{eq:deterministic-finite}
\end{equation}
\end{proposition}

\proofsketch
Kink placement at $q_0$ gives
$\widetilde q^{\mathrm{LL}}(\lambda)\le q_0<q_1$ and
$\widetilde\mu\le\mu_{q_0}$.  For $q_1\le0$, the map
\[
 \mu\longmapsto \mu\ell_0'(\lambda+\mu q_1)
\]
is nondecreasing on the nonbinding range: increasing $\mu$ directly increases
the first factor and weakly lowers the argument of the decreasing derivative.
Combine this observation with Lemma~\ref{lem:curvature-split} at $q_1$ and
$c''(a)\ge\underline c''$.

\section{High feasible reservation utility}

\begin{theorem}[Unified asymptotic concavity criterion]
\label{thm:unified}
Suppose Assumption~\ref{ass:link-concavity} and \eqref{eq:cost-curv} hold.  Fix
$q_0<q_1\le0$ with $q_0,q_1\in\cQ_{\mathrm{safe}}$.  Assume
\[
 \Delta u\,\Capacity(q_0)>c'(a_0),
\]
and
\begin{equation}
 \limsup_{\lambda\to\infty}
 \mu_{q_0}(\lambda)
 \ell_0'(\lambda+\mu_{q_0}(\lambda)q_1)\Curv(q_1)
 <\underline c''.
 \tag{DC}\label{eq:DC}
\end{equation}
Then there is $\lambda_0$ such that, for every $\lambda\ge\lambda_0$,
\[
 U_{aa}(v_\lambda,a)\le0
 \quad\text{for every }a\in\cA.
\]
Consequently, every corresponding relaxed solution satisfies global incentive
compatibility.
\end{theorem}

\proofsketch
The strict inequality in \eqref{eq:DC} implies $\Curv(q_1)<\infty$ and,
together with Proposition~\ref{prop:deterministic-criterion}, gives the conclusion.  Since LIC holds at the interior action
$a_0$, concavity of $a\mapsto U(v_\lambda,a)$ makes $a_0$ a global maximizer.

\subsection{Strong flattening}

\begin{definition}[Strong flattening]
Utility satisfies strong flattening if
\begin{equation}
 \frac{[u'(W)]^2}{-u''(W)}\longrightarrow0
 \quad\text{as }W\to\infty.
 \tag{SF-$u$}\label{eq:SF-u}
\end{equation}
Equivalently,
\begin{equation}
 z\ell_0'(z)\longrightarrow0
 \quad\text{as }z\to\infty.
 \tag{SF}\label{eq:SF}
\end{equation}
Indeed, when $v=u(W)$ and $z=k'(v)=1/u'(W)$,
\[
 z\ell_0'(z)=\frac{[u'(W)]^2}{-u''(W)}.
\]
\end{definition}

\begin{theorem}[Strong-flattening utilities]
\label{thm:strong-flattening}
Suppose Assumption~\ref{ass:link-concavity}, \eqref{eq:cost-curv}, and strong
flattening hold.  Suppose
\[
 \Delta u\,\Capacity_{\mathrm{safe}}>c'(a_0),
 \qquad
 \Curv_{\mathrm{safe}}<\infty.
\]
Then the LIC-path contracts induce an agent objective concave in action for all
sufficiently large $\lambda$.
\end{theorem}

\proofsketch
The aggregate conditions provide $q_0<q_1$ in $\cQ_{\mathrm{safe}}$ such that
$\Delta u\,\Capacity(q_0)>c'(a_0)$ and $\Curv(q_1)<\infty$.  The trial
multiplier is $\mu_{q_0}(\lambda)=O(\lambda)$, while
\[
 z_{01}(\lambda):=\lambda+\mu_{q_0}(\lambda)q_1
 \sim \lambda\left(1-\frac{q_1}{q_0}\right).
\]
Because $q_0<q_1$, the coefficient is positive.  Therefore
\[
 \mu_{q_0}(\lambda)\ell_0'(z_{01}(\lambda))\Curv(q_1)\longrightarrow0
\]
by $z\ell_0'(z)\to0$.  Apply Theorem~\ref{thm:unified}.

\subsection{The logarithmic borderline}

For log utility, $\ell_0'(z)=1/z$, so $z\ell_0'(z)=1$ and strong flattening fails.

\begin{theorem}[Log utility and safe score width]
\label{thm:log-width}
Suppose $u(W)=\log W$, Assumption~\ref{ass:link-concavity}, and
\eqref{eq:cost-curv}.  Suppose $\cQ_{\mathrm{safe}}$ is nonempty and
\begin{equation}
 \Curv_{\mathrm{safe}}
 <\underline c''\,\Width_{\mathrm{safe}}.
 \tag{LogWidth}\label{eq:log-width}
\end{equation}
If $\underline s=-\infty$, then $\Width_{\mathrm{safe}}=\infty$ and the right
side is interpreted as infinity.
Then the LIC-path contracts induce an agent objective concave in action for all
sufficiently large $\lambda$.
\end{theorem}

\proofsketch
The strict endpoint inequality provides
$\underline s<q_0<q_1\le0$ with $q_0,q_1\in\cQ_{\mathrm{safe}}$ such that
$\Capacity(q_0)>0$ and
\[
 \frac{\Curv(q_1)}{q_1-q_0}<\underline c''.
\]
Log utility is unbounded, so the first condition gives capacity.  Moreover,
\[
 \mu_{q_0}(\lambda)
 \ell_0'(\lambda+\mu_{q_0}(\lambda)q_1)
 =\frac{\mu_{q_0}(\lambda)}
 {\lambda+\mu_{q_0}(\lambda)q_1}
 \longrightarrow\frac1{q_1-q_0}.
\]
Thus \eqref{eq:DC} holds.

\subsection{Passage to reservation utility and the FOA}

\begin{proposition}[Concavity implies global incentive compatibility]
\label{prop:concavity-GIC}
Let $v$ satisfy LIC at the interior action $a_0$.  If
$a\mapsto U(v,a)$ is concave on $\cA$, then $a_0$ is a global maximizer and $v$
satisfies global incentive compatibility.
\end{proposition}

\begin{theorem}[FOA at high feasible reservation utility]
\label{thm:high-reservation}
Suppose Proposition~\ref{prop:relaxed-solution} applies.  Suppose, in addition,
that either the hypotheses of Theorem~\ref{thm:strong-flattening} or those of
Theorem~\ref{thm:log-width} hold.  Then there exists
$U^*<U_R$ such that the first-order approach is valid for every
\[
 \bar U\in[U^*,U_R).
\]
If $\bar u=\infty$, then $U_R=\infty$ and this is the usual sufficiently-high
reservation-utility statement.  If $\bar u<\infty$, it means reservation utility
sufficiently close to the relaxed upper frontier from below.
\end{theorem}

\proofsketch
The relevant concavity theorem supplies $\lambda_0$.  Proposition~\ref{prop:relaxed-solution}
gives $\lambda^*(\bar U)\to\infty$ as $\bar U\uparrow U_R$.  Choose $U^*$ so
that $\lambda^*(\bar U)\ge\lambda_0$ above $U^*$, and apply
Proposition~\ref{prop:concavity-GIC}.

\section{Utility corollaries}

The following statements hold the distribution, intended action, effort cost,
and utility-shape parameter fixed.  They maintain the regularity conditions for
the relaxed solution and uniform cost curvature~\eqref{eq:cost-curv}.

\begin{corollary}[CARA at sufficiently low initial wealth]
\label{cor:cara-low-wealth}
Suppose $\Capacity_{\mathrm{safe}}>0$ and
$\Curv_{\mathrm{safe}}<\infty$.  Fix $\alpha>0$ and let
$u(W)=-e^{-\alpha W}/\alpha$.  For all sufficiently low initial wealth $w_0$,
the first-order approach is valid at all sufficiently high feasible reservation
utilities.
\end{corollary}

\proofsketch
CARA utility satisfies strong flattening, and its utility spread diverges as
$w_0\to-\infty$.  Apply Theorems~\ref{thm:strong-flattening}
and~\ref{thm:high-reservation}.

\begin{corollary}[CRRA above one at sufficiently low initial wealth]
\label{cor:crra-low-wealth}
Suppose $\Capacity_{\mathrm{safe}}>0$ and
$\Curv_{\mathrm{safe}}<\infty$.  Fix $\gamma>1$ and let
$u(W)=W^{1-\gamma}/(1-\gamma)$.  For all sufficiently low positive initial
wealth $w_0$, the first-order approach is valid at all sufficiently high feasible
reservation utilities.
\end{corollary}

\proofsketch
CRRA utility with $\gamma>1$ satisfies strong flattening, and its utility spread
diverges as $w_0\downarrow0$.  Apply Theorems~\ref{thm:strong-flattening}
and~\ref{thm:high-reservation}.

\begin{corollary}[Log utility with an unbounded left score]
\label{cor:log-unbounded-left}
Suppose $w_0>0$, $u(W)=\log W$, $\underline s=-\infty$, and
$\Curv_{\mathrm{safe}}<\infty$.  Then there is $U^*<\infty$ such that the
first-order approach is valid for every reservation utility $\bar U\ge U^*$.
\end{corollary}

\proofsketch
The safe score width is infinite, so condition~\eqref{eq:log-width} holds.  Log
utility is unbounded above, so $U_R=\infty$.  Apply
Theorem~\ref{thm:high-reservation}.

\section{Distribution verifications}

\subsection{Gaussian location}

Suppose $Y\sim N(a,\sigma^2)$.  Then
\[
 s(y)=\frac{y-a_0}{\sigma^2},
 \qquad
 f_{aa}(y\mid a)=f(y\mid a)
 \left[\frac{(y-a)^2}{\sigma^4}-\frac1{\sigma^2}\right].
\]
If $\cA=[\underline a,\bar a]$, the right edge of the uniformly safe low-score
region is
\[
 \bar q_{\mathrm{safe}}
 =\min\left\{0,\frac{\underline a-a_0-\sigma}{\sigma^2}\right\}.
\]
Under $a_0$, $s(Y)\sim N(0,1/\sigma^2)$, and
\[
 \Capacity(q)=\frac1\sigma\varphi(\sigma q).
\]
Paid-region curvature is finite on compact action sets, so
$\Curv_{\mathrm{safe}}<\infty$.  Since the score is unbounded below, the log-width
condition follows.  For CARA and CRRA with $\gamma>1$, safe incentive capacity
remains a substantive condition.

\subsection{Poisson}

Suppose $Y\sim\operatorname{Pois}(a)$ and $\underline a:=\min\cA>0$.  Then
\[
 s(y)=\frac{y-a_0}{a_0},
 \qquad
 p_{aa}(y\mid a)=p(y\mid a)\frac{(y-a)^2-y}{a^2}.
\]
A count $y$ is uniformly low-tail safe if
\[
 y+\sqrt y\le\underline a.
\]
For a safe integer cutoff $K<a_0$, let $q_K:=(K-a_0)/a_0$ be the
corresponding score cutoff.  Then
\[
 \Capacity(q_K)
 =-\sum_{y\le K}\frac{y-a_0}{a_0}p(y\mid a_0)
 =\Pr_{a_0}(Y=K).
\]
Paid-region curvature is finite on compact action sets bounded away from zero.
Because score values are discrete, formal applications should compute interior
safe capacity rather than mechanically using capacity at the rightmost safe
atom.

\subsection{Exponential output with mean action}

Suppose output has density
\[
 f(y\mid a)=\frac1a e^{-y/a}\ind\{y\ge0\},
 \qquad \underline a:=\min\cA>0.
\]
Then
\[
 s(y)=\frac{y-a_0}{a_0^2},
 \qquad
 f_{aa}(y\mid a)=f(y\mid a)
 \frac{y^2-4ay+2a^2}{a^4}.
\]
The uniformly safe low-output region is
\[
 y\le(2-\sqrt2)\underline a,
\]
so
\[
 \bar q_{\mathrm{safe}}
 =\min\left\{0,
 \frac{(2-\sqrt2)\underline a-a_0}{a_0^2}
 \right\}.
\]
For a nonnegative output cutoff $r=a_0+a_0^2q$ and $b=r/a_0$,
\[
 \Capacity(q)=\frac{be^{-b}}{a_0}.
\]
Again, $\Curv(q)<\infty$ on compact action sets bounded away from zero.

\section{Appendix-only asymptotic-affinity extension}

This section uses a different mechanism from safe-tail flattening.  It is intended
for an appendix, with at most a brief main-text reference.

\subsection{Affine-score curvature and primitive moments}

\begin{assumption}[Affine-score curvature]
\label{ass:affine-score}
For every $a\in\cA$,
\begin{equation}
 \int s(y)f_{aa}(y\mid a)\dd(y)\le0.
 \tag{Affine}\label{eq:affine-score}
\end{equation}
Equivalently, $a\mapsto\E_a[s(Y)]$ is concave.
\end{assumption}

\begin{assumption}[Moments for asymptotic affinity]
\label{ass:affinity-moments}
For some $\eta>0$,
\begin{align*}
 0<\sigma_s^2&:=\int s(y)^2f_0(y)\dd(y)<\infty,\\
 \int |s(y)|^{2+\eta}f_0(y)\dd(y)&<\infty,\\
 \sup_{a\in\cA}\int
 (1+|s(y)|^{2+\eta})|f_{aa}(y\mid a)|\dd(y)&<\infty.
\end{align*}
\end{assumption}

\subsection{CRRA with one-half below gamma below one}

Suppose
\[
 \frac12<\gamma<1,
 \qquad
 p:=\frac{1-\gamma}{\gamma}\in(0,1).
\]
On the nonbinding region,
\[
 \ell_0(z)=\frac{z^p}{1-\gamma},
 \qquad
 \ell_0'(z)=\frac1\gamma z^{p-1}.
\]
Thus $\ell_0'(z)\to0$ but $z\ell_0'(z)\to\infty$, so strong flattening fails.

\begin{lemma}[LIC multiplier asymptotics]
\label{lem:affinity-multiplier}
Under Assumption~\ref{ass:affinity-moments}, for CRRA with
$\gamma\in(1/2,1)$,
\[
 \widetilde\mu(\lambda)\ell_0'(\lambda)
 \longrightarrow
 \theta^*:=\frac{c'(a_0)}{\sigma_s^2},
 \qquad
 \frac{\widetilde\mu(\lambda)}{\lambda}\longrightarrow0.
\]
\end{lemma}

\proofsketch
For bounded $\theta$, set $\mu_\lambda(\theta)=\theta/\ell_0'(\lambda)$.  Since
$\lambda\ell_0'(\lambda)\to\infty$, we have $\mu_\lambda(\theta)/\lambda\to0$.
A central-region Taylor expansion and the $(2+\eta)$ score moment give
\[
 \ell(\lambda+\mu_\lambda(\theta)s)
 =\ell(\lambda)+\theta s+o(1)
\]
in the weighted space
\[
 L^1\bigl(|s(y)|f_0(y)\dd(y)\bigr).
\]
The constant vanishes after multiplication by $s$ and integration.  Hence
$I(\lambda,\mu_\lambda(\theta))\to\theta\sigma_s^2$.
Monotonicity of $I$ pins down the LIC solution.  The second limit follows by
dividing the first by $\lambda\ell_0'(\lambda)$.

\begin{lemma}[Uniform residual]
\label{lem:uniform-residual}
Under Assumption~\ref{ass:affinity-moments},
\begin{equation}
\begin{aligned}
 &\sup_{a\in\cA}\Bigg|
 \int\left[
  \ell(\lambda+\widetilde\mu(\lambda)s(y))-\ell(\lambda)
  -\widetilde\mu(\lambda)\ell_0'(\lambda)s(y)
 \right]
 \\[-2pt]
 &\hspace{5cm}{}
 \times f_{aa}(y\mid a)\dd(y)
 \Bigg|\longrightarrow0.
\end{aligned}
\tag{R}\label{eq:residual}
\end{equation}
\end{lemma}

\proofsketch
Split at $|s|\le \lambda/(2\widetilde\mu)$.  On the central region, Taylor's theorem
bounds the residual by a constant times
$\widetilde\mu^2|\ell_0''(\lambda/2)|s^2$, whose coefficient is $O(\lambda^{-p})$.
On the tail, use the sublinear growth $|\ell(z)|\le C(1+|z|^p)$ and the uniform
$(2+\eta)$ weighted moment of $f_{aa}$.  Lemma~\ref{lem:affinity-multiplier}
gives the rates needed for both pieces to vanish uniformly in $a$.

\begin{theorem}[Asymptotic affinity]
\label{thm:asymptotic-affinity}
Suppose \eqref{eq:cost-curv}, Assumptions~\ref{ass:affine-score} and
\ref{ass:affinity-moments}, and CRRA utility with $\gamma\in(1/2,1)$.  Then
\[
 \sup_{a\in\cA}
 U_{aa}(v_\lambda,a)
 \le-\underline c''+o(1).
\]
In particular, the agent's objective is concave for all sufficiently large
$\lambda$, and the FOA is valid at all sufficiently high feasible reservation
utilities.
\end{theorem}

\proofsketch
Use $\int f_{aa}\dd=0$ to remove $\ell(\lambda)$ and expand
\begin{align*}
 \int \ell(\lambda+\widetilde\mu s)f_{aa}(y\mid a)\dd(y)
 &=\widetilde\mu\ell_0'(\lambda)
   \int s(y)f_{aa}(y\mid a)\dd(y)+o(1).
\end{align*}
The remainder is uniform by Lemma~\ref{lem:uniform-residual}; the leading term is
nonpositive by affine-score curvature.  Subtract $c''(a)\ge\underline c''$.

\begin{corollary}[Gaussian CRRA extension]
For a Gaussian location family on a compact action set, all moments in
Assumption~\ref{ass:affinity-moments} hold and
\[
 \int s(y)f_{aa}(y\mid a)\dd(y)
 =\frac{d^2}{da^2}\E_a[s(Y)]=0.
\]
Consequently, Gaussian output with CRRA utility is covered by strong flattening
for $\gamma>1$, by the log theorem for $\gamma=1$, and by asymptotic affinity for
$\gamma\in(1/2,1)$, subject in each branch to the assumptions of the corresponding
relaxed-solution result.
\end{corollary}

No general claim is made here for $\gamma\le1/2$.

\section{Result map for the revised paper}

The recommended manuscript order is:
\begin{enumerate}
 \item model, score, links, and canonical contracts;
 \item generalized relaxed-solution Proposition~\ref{prop:relaxed-solution};
 \item tail-safe cutoffs and the finite-$\lambda$ kink and curvature lemmas;
 \item unified criterion, followed by strong flattening and log as transparent
 verification results;
 \item high-feasible-reservation Theorem~\ref{thm:high-reservation};
 \item utility and distribution applications;
 \item asymptotic affinity only in the appendix.
\end{enumerate}

The main conceptual decomposition is
\[
 \boxed{\text{the kink lies in a safe low-score tail}}
 \quad+\quad
 \boxed{\text{paid-region positive curvature is dominated by }c''}.
\]
The first box is supplied by incentive capacity and kink placement.  The second
is supplied either by strong flattening, by sufficient safe score width for log
utility, or---in the appendix---by asymptotic affinity.

\end{document}
